% Figure: convergence of the shifted QR algorithm, measured by the
% magnitude of the trailing subdiagonal entry against iteration count,
% on a log-y axis. Unshifted QR converges linearly; the Wilkinson
% shift gives the near-cubic plunge that deflates one eigenvalue in a
% few steps.
\begin{figure}[htbp]
\centering
\begin{tikzpicture}
\begin{semilogyaxis}[
  width=13cm,
  height=8cm,
  xlabel={QR iteration $t$},
  ylabel={trailing subdiagonal $|T_{n,n-1}^{(t)}|$},
  xmin=0, xmax=12,
  ymin=1e-16, ymax=1,
  ytick={1,1e-3,1e-6,1e-9,1e-12,1e-15},
  legend pos=south west,
  legend style={font=\footnotesize},
  tick label style={font=\footnotesize},
  label style={font=\small},
  title={Subdiagonal decay: unshifted vs.\ shifted QR},
  title style={font=\bfseries},
]
% Unshifted QR: linear convergence at rate |lambda_n / lambda_{n-1}| ~ 0.6.
\addplot[engblue, very thick, mark=*, mark size=1.4pt] coordinates {
  (0, 0.6) (1, 0.36) (2, 0.216) (3, 0.13) (4, 0.078)
  (5, 0.047) (6, 0.028) (7, 0.017) (8, 0.010) (9, 0.0061)
  (10, 0.0037) (11, 0.0022) (12, 0.0013)
};
\addlegendentry{unshifted (linear)}
% Wilkinson-shifted QR: near-cubic, squaring/cubing the error each step.
\addplot[engred, very thick, mark=square*, mark size=1.4pt] coordinates {
  (0, 0.6) (1, 0.12) (2, 6e-3) (3, 8e-6) (4, 2e-14) (5, 1e-16)
};
\addlegendentry{Wilkinson shift (near cubic)}
\end{semilogyaxis}
\end{tikzpicture}
\caption[Subdiagonal decay in the shifted QR algorithm]{The QR
algorithm computes eigenvalues by driving the subdiagonal entries of a
Hessenberg matrix to zero; once the trailing entry $T_{n,n-1}$ is
negligible, the bottom-right corner holds a converged eigenvalue and
the problem deflates to size $n-1$. Unshifted, the trailing entry
shrinks geometrically at the ratio of the two smallest eigenvalue
magnitudes (blue), so a tight cluster converges slowly. The Wilkinson
shift, which subtracts an estimate of the target eigenvalue before each
factor-and-reverse step, turns the linear decay into near-cubic
plunge (red): the error is roughly cubed each iteration, and one
eigenvalue deflates in three to five steps regardless of the cluster.
Shifting is the single change that makes the QR algorithm a practical
$\oom{n^{3}}$ eigensolver rather than a textbook curiosity, and it is
what every LAPACK dense eigenroutine runs underneath.}
\label{fig:vol02:ch10:qr-convergence}
\end{figure}
